\section{Формальная постановка задачи и алгоритмы}

Настоящий раздел вводит три объекта, которыми оперирует работа, и выводит
из них задачи визуализации. Каждая задача описывается входом, выходом и
алгоритмом, причём вход задаётся только в терминах ранее определённых
объектов: величин, возникающих вне этой цепочки, в постановках нет.

\subsection{Исходные объекты}

\subsubsection{Дерево игры}

Игра задаётся набором компонент $G = \langle R, S, s_0, A, noop,
terminal(\cdot), legal_i(\cdot), next(\cdot), payoff(\cdot) \rangle$.
Порождаемое ею дерево игры $T_{game} = (V_{game}, E_{game})$ определяется
так. Вершина~--- это \emph{история}, то есть последовательность совместных
действий $h = \langle a_1, \dots, a_k \rangle$, выполнимая из начального
состояния; корень~--- пустая история $\varepsilon$. Каждой истории отвечает
состояние игры

\begin{equation}
state(\varepsilon) = s_0, \qquad
state(h \cdot a) = next(state(h), a).
\end{equation}

Дуги соединяют историю с её продолжениями допустимым совместным действием:

\begin{equation}
E_{game} = \bigl\{ (h,\; h \cdot a) \;\big|\; h \in V_{game},\;
a \in legal(state(h)) \bigr\},
\end{equation}

а листьями служат истории с $terminal(state(h)) = \mathit{true}$.

Вершиной объявлена история, а не состояние, намеренно. Одно и то же
состояние достижимо разными последовательностями ходов (транспозиция), и
отождествление таких вершин превратило бы дерево в ориентированный
ациклический граф. Отношение эквивалентности состояний
$h_1 \sim h_2 \Leftrightarrow state(h_1) = state(h_2)$ сохраняется отдельно
и используется как дополнительный слой, не меняющий древесной структуры.

Дерево игры \emph{дано}: оно определено правилами и не строится ни
алгоритмом поиска, ни визуализацией. Его мощность для рассматриваемых игр
исключает перечисление вершин.

\subsubsection{Дерево поиска}

Алгоритм поиска строит конечное поддерево дерева игры, снабжённое
статистикой. Источником сведений о построении служит журнал итераций~---
последовательность событий $\sigma = \langle e_1, e_2, \dots \rangle$ трёх
видов:

\begin{itemize}
  \item $\mathrm{TurnStart}(h)$~--- история $h$ становится корнем поиска на
очередном ходе партии, остальное дерево отсекается;
  \item $\mathrm{Expand}(i, h, C)$~--- на итерации $i$ у вершины $h$
появляются потомки $C$, отвечающие допустимым продолжениям;
  \item $\mathrm{Iter}(i, l, p, s, d)$~--- итерация $i$ выбрала лист $l$,
выполнила симуляцию из вершины $p$, получила платежи $s$ и глубину
доигрывания $d$.
\end{itemize}

Дерево поиска на момент $t$ (числом выполненных итераций) есть

\begin{equation}
T_{search}(t) = \bigl( V_t,\; E_t,\; N_t,\; Q_t \bigr), \qquad
V_t \subseteq V_{game},
\end{equation}

где $V_t$~--- вершины, порождённые событиями расширения до момента $t$,
$E_t$~--- индуцированные на них дуги дерева игры, $N_t\colon V_t \to
\mathbb{N}$~--- число посещений, $Q_t\colon V_t \to \mathbb{R} \cup
\{\bot\}$~--- средняя накопленная оценка ($\bot$, если оценки нет).

Пути этапов выбора и обратного распространения в журнале не хранятся: оба
проходят по единственному пути от корня к вершине симуляции и
восстанавливаются по цепочке предков. Это свойство выводимости определяет
состав журнала и используется далее при реконструкции.

Поиск по дереву методом Монте-Карло рассматривается как частный случай:
изложение опирается только на перечисленные события, поэтому любая
вариация метода, порождающая такой журнал, укладывается в ту же модель.

\subsubsection{Геометрическое дерево}

Геометрическое дерево~--- это размещение вершин дерева поиска на плоскости:

\begin{equation}
T_{geom}(t) = \bigl( V_t,\; x_t,\; y \bigr), \qquad
x_t\colon V_t \to \{[a, b] \subseteq [0, 1]\}, \quad y(v) = depth(v).
\end{equation}

Вершине отвечает прямоугольник: по горизонтали~--- интервал $x_t(v) = [x_0(v),
x_1(v)]$ нормированной оси, по вертикали~--- ярус номера $y(v)$ постоянной
высоты. Ширина интервала $w(v) = x_1(v) - x_0(v)$.

Геометрия подчинена двум требованиям:

\begin{equation}
x_t(c) \subseteq x_t(p) \quad \text{для всякой дуги } (p, c) \in E_t,
\label{eq:nesting}
\end{equation}
\begin{equation}
\sum_{c \in children(p)} w(c) = w(p).
\label{eq:partition}
\end{equation}

Вложенность~(\ref{eq:nesting}) кодирует отношение <<предок~--- потомок>>
положением прямоугольников, поэтому дуги не изображаются вовсе; разбиение
без остатка~(\ref{eq:partition}) означает, что потомки делят полосу предка
полностью. Именно отказ от явных дуг отличает это представление от
узло-связевого и освобождает экранное пространство.

Отображение геометрического дерева на экран задаётся камерой~--- видимым
доменом $D = [d_0, d_1] \subseteq [0,1]$ оси долей. При ширине окна $W$
пикселей плотность отображения и экранная ширина вершины равны

\begin{equation}
\rho = \frac{W}{d_1 - d_0}, \qquad w_{px}(v) = w(v) \cdot \rho.
\end{equation}

\subsection{Общая задача и алгоритм верхнего уровня}

\paragraph{Задача В (визуализировать состояние поиска).}
\emph{Вход:} журнал $\sigma$; момент $t$; камера $(D, W)$; бюджет
изображаемых вершин $B$.
\emph{Выход:} множество изображаемых фигур~--- прямоугольников вершин
$R$ и агрегатов скрытых частей $H$, покрывающих полосу корня без пустот.

Задача В решается последовательным применением четырёх подзадач; выход
каждой служит входом следующей, и ни одна величина не вводится извне.

\begin{quote}
\textbf{Алгоритм~1. Визуализация состояния поиска}\\
\emph{Вход:} $\sigma$, $t$, $D$, $W$, $B$. \emph{Выход:} $(R, H)$.
\begin{enumerate}
  \item $T \leftarrow \textsc{Реконструировать}(\sigma, t)$
        \hfill (задача Р, выход $T_{search}(t)$)
  \item $g \leftarrow \textsc{ПостроитьГеометрию}(T, \sigma, t)$
        \hfill (задача Г, выход $T_{geom}(t)$)
  \item $\rho \leftarrow W / (d_1 - d_0)$
  \item $R \leftarrow \textsc{ОтобратьВидимые}(g, D, \rho, B)$
        \hfill (задача О, выход $R \subseteq V_t$)
  \item $H \leftarrow \textsc{СвернутьСкрытое}(g, R)$
        \hfill (задача С, выход агрегатов)
  \item \textbf{вернуть} $(R, H)$
\end{enumerate}
\end{quote}

\subsection{Подзадачи}

\subsubsection{Задача Р: реконструировать состояние поиска}

\emph{Вход:} журнал $\sigma$, момент $t$.
\emph{Выход:} дерево поиска $T_{search}(t)$.

Событие расширения подвешивает к вершине перечисленных потомков; событие
итерации увеличивает на единицу счётчик посещений каждой вершины на пути от
вершины симуляции к корню и обновляет там статистику выбранного действия;
событие границы хода делает указанную вершину корнем и отсекает остальное.

\begin{quote}
\textbf{Алгоритм~2. Реконструкция}\\
\emph{Вход:} $\sigma$, $t$. \emph{Выход:} $T_{search}(t)$.
\begin{enumerate}
  \item $T \leftarrow$ пустое дерево
  \item \textbf{для} каждого события $e \in \sigma$ по порядку, пока не
        достигнут момент $t$:
  \begin{enumerate}
    \item \textbf{если} $e = \mathrm{TurnStart}(h)$: сделать $h$ корнем $T$,
          отсечь вершины вне поддерева $h$
    \item \textbf{если} $e = \mathrm{Expand}(i, h, C)$: подвесить $C$ к $h$
          с нулевыми счётчиками
    \item \textbf{если} $e = \mathrm{Iter}(i, l, p, s, d)$: для каждой
          вершины $v$ на пути от $p$ к корню увеличить $N(v)$ на единицу и
          учесть $s$ в статистике действия, ведущего в потомка на этом пути
  \end{enumerate}
  \item \textbf{вернуть} $T$
\end{enumerate}
\end{quote}

Трудоёмкость линейна по числу прочитанных событий. Перемотка вперёд
продолжает чтение с текущей позиции; перемотка назад требует чтения с
начала и ускоряется контрольными точками~--- сохранёнными копиями
состояния, от ближайшей из которых чтение продолжается.

\subsubsection{Задача Г: построить опорную геометрию}

\emph{Вход:} дерево поиска $T_{search}(t)$; журнал $\sigma$ и момент $t$~---
для выбора опорного состояния.
\emph{Выход:} геометрическое дерево $T_{geom}(t)$, то есть интервалы $x_t$
и ярусы $y$, удовлетворяющие~(\ref{eq:nesting}) и~(\ref{eq:partition}).

Доля потомка считается среди соседей по общему предку, а не относительно
корня: при переиспользовании дерева между ходами счётчики предков
накапливаются за всю партию, и абсолютные значения несопоставимы между
ярусами.

\begin{quote}
\textbf{Алгоритм~3. Назначение интервалов}\\
\emph{Вход:} $T_{search}$, опорное состояние $T^{*}$. \emph{Выход:} $x$, $y$.
\begin{enumerate}
  \item корню назначить $x(r) = [0, 1]$, $y(r) = 0$
  \item обходом в глубину по $T^{*}$ для каждой вершины $p$ с потомками
        $c_1, \dots, c_n$ в естественном порядке ходов:
  \begin{enumerate}
    \item $\Sigma \leftarrow \sum_{k} N(c_k)$
    \item $\varphi(c_k) \leftarrow N(c_k) / \Sigma$, если $\Sigma > 0$,
          иначе $\varphi(c_k) \leftarrow 1/n$
    \item $\delta \leftarrow \min(0{,}02;\; 0{,}5/n)$;\quad
          $s(c_k) \leftarrow \delta + (1 - n\delta)\,\varphi(c_k)$
    \item разделить $x(p)$ между потомками слева направо по долям $s(c_k)$,
          $y(c_k) \leftarrow y(p) + 1$
  \end{enumerate}
  \item вершинам $T_{search} \setminus T^{*}$ интервалы не назначаются
  \item \textbf{вернуть} $x$, $y$
\end{enumerate}
\end{quote}

Нижний порог $\delta$ отделяет два канала: ширина отвечает на вопрос,
рассматривался ли ход, а доля $\varphi$ без порога~--- насколько он изучен и
служит основой цветового кодирования.

Опорное состояние $T^{*}$ выбирается в зависимости от режима просмотра и
определяет устойчивость геометрии (задача У).

\subsubsection{Задача О: отобрать изображаемые вершины}

\emph{Вход:} $T_{geom}(t)$; домен $D$; плотность $\rho$; бюджет $B$.
\emph{Выход:} связное поддерево $R \subseteq V_t$, $|R| \le B$.

Вершина допустима, если её интервал пересекает домен и её экранная ширина
достигает перцептивного минимума $s_{min}$:

\begin{equation}
adm(v) \;\Leftrightarrow\; x(v) \cap D \ne \varnothing \;\wedge\;
w(v) \cdot \rho \ge s_{min}.
\label{eq:adm}
\end{equation}

\begin{quote}
\textbf{Алгоритм~4. Отбор по экранному размеру с бюджетом}\\
\emph{Вход:} $T_{geom}$, $D$, $\rho$, $B$. \emph{Выход:} $R$.
\begin{enumerate}
  \item $R \leftarrow \varnothing$; очередь с приоритетом по убыванию
        ширины содержит корень, если он допустим
  \item \textbf{пока} очередь не пуста и $|R| < B$:
  \begin{enumerate}
    \item извлечь вершину $v$ наибольшей ширины, $R \leftarrow R \cup \{v\}$
    \item поместить в очередь всех потомков $v$, удовлетворяющих
          условию~(\ref{eq:adm})
  \end{enumerate}
  \item \textbf{вернуть} $R$
\end{enumerate}
\end{quote}

Связность выхода обеспечена построением: из вложенности~(\ref{eq:nesting})
следует $w(p) \ge w(c)$, поэтому допустимость потомка влечёт допустимость
предка, а потомок попадает в очередь лишь после включения предка.
Монотонность по масштабу также следует из~(\ref{eq:adm}): при росте $\rho$
множество допустимых вершин только расширяется, а вновь допустимые заведомо
у́же ранее включённых и не вытесняют их из бюджета.

\subsubsection{Задача С: свернуть скрытую часть}

\emph{Вход:} $T_{geom}(t)$; отобранное множество $R$.
\emph{Выход:} множество агрегатов $H$, каждый со своей сводкой.

Участки полосы вершины $v \in R$, не покрытые её изображаемыми потомками,
образуют скрытую массу. Агрегат строится не на всю вершину, а на каждую
серию подряд идущих неизображаемых потомков: у вершины, рассечённой
изображаемыми потомками, таких участков несколько.

\begin{quote}
\textbf{Алгоритм~5. Свёртка скрытого}\\
\emph{Вход:} $T_{geom}$, $R$. \emph{Выход:} $H$.
\begin{enumerate}
  \item \textbf{для} каждой $v \in R$ и каждой максимальной серии подряд
        идущих потомков $c_j, \dots, c_m \notin R$:
  \begin{enumerate}
    \item интервал агрегата $[x_0(c_j),\, x_1(c_m)]$
    \item сводка: число свёрнутых ходов $m - j + 1$; число вершин в их
          поддеревьях; суммарные посещения и их доля среди всех потомков
          $v$; глубина глубочайшего потомка; диапазон оценок
  \end{enumerate}
  \item \textbf{вернуть} $H$
\end{enumerate}
\end{quote}

Сводка обязательна: агрегат, сообщающий лишь факт сокрытия, не позволяет
судить о свёрнутом, а значит не сохраняет достоверности представления. Из
разбиения~(\ref{eq:partition}) следует инвариант сохранения массы: на
каждом ярусе сумма ширин изображаемых вершин и агрегатов равна ширине
полосы предка, то есть пользователь всегда видит всю статистику поиска,
меняется лишь её гранулярность.

\subsubsection{Задача У: обеспечить устойчивость при росте}

\emph{Вход:} последовательность моментов $t_1 < t_2 < \dots$ внутри одного
хода; журнал $\sigma$.
\emph{Выход:} выбор опорного состояния $T^{*}$ для задачи Г, при котором
положение вершины не зависит от момента $t$.

Если опорным состоянием брать текущее дерево $T_{search}(t)$, доли
изменяются с каждой итерацией и вершины смещаются от кадра к кадру. Для
записанного хода доступно его итоговое состояние $T_{search}(T_h)$, и
поскольку внутри хода вершины только добавляются, справедливо
$V_t \subseteq V_{T_h}$ при $t \le T_h$. Отсюда опорное состояние

\begin{equation}
T^{*} = T_{search}(T_h) \quad\text{для всех } t \in (T_{h-1}, T_h],
\end{equation}

и геометрия неподвижна на протяжении всего хода: вершина появляется сразу
на своём окончательном месте, а изменяется только цвет, несущий текущую
долю посещений. При наблюдении в реальном времени итоговое состояние хода
ещё не существует, поэтому опорным служит текущее дерево, и устойчивость в
этом режиме не гарантируется~--- режимы различаются именно выбором $T^{*}$.

\subsection{Связность постановок}

Согласованность входов и выходов проверяется по таблице~\ref{tab:io}: всякий
вход либо принадлежит исходным данным (журнал, камера, бюджет), либо
является выходом ранее решённой задачи.

\begin{gosttable}{tab:io}{Входы и выходы задач}
\begin{tabularx}{\linewidth}{>{\raggedright\arraybackslash}p{0.10\linewidth}>{\raggedright\arraybackslash}X>{\raggedright\arraybackslash}p{0.26\linewidth}}
\hline
Задача & Вход & Выход \\
\hline
Р & журнал $\sigma$, момент $t$ & дерево поиска $T_{search}(t)$ \\
\hline
У & журнал $\sigma$, границы хода & опорное состояние $T^{*}$ \\
\hline
Г & $T_{search}(t)$, $T^{*}$ & геометрическое дерево $T_{geom}(t)$ \\
\hline
О & $T_{geom}(t)$, камера $(D, W)$, бюджет $B$ & изображаемое множество $R$ \\
\hline
С & $T_{geom}(t)$, множество $R$ & агрегаты $H$ со сводками \\
\hline
\end{tabularx}
\end{gosttable}

Исходными данными служат только журнал, параметры камеры и бюджет. Дерево
игры при этом не вычисляется: оно присутствует в постановках как
вместилище, подмножеством которого является дерево поиска.
